\section{What synthetic prediction can teach}
\label{sec:information}

\subsection{Determinism, rule uncertainty and finite entropy}

Condition on a fixed topology and rule. The complete trajectory is a
deterministic function of $X^0$, so, in bits,
\begin{equation}
 H(X^{0:T-1}\mid F,G)=H(X^0\mid F,G)
 \leq n\log_2q.
 \label{eq:knownruleentropy}
\end{equation}
The equality includes $X^0$ in the observed trajectory. It does not mean
that a computationally bounded predictor can cheaply recover all later
states, or that a short-window predictor sees $X^0$.

Continuous neural weights do not have a discrete entropy measured by their
parameter count. Instead, consider the \emph{effective deterministic local
rule} $F:\mathcal A^{d+1}\to\mathcal A$.
There are at most $q^{q^{d+1}}$ such functions, including functions not
representable by the chosen MLP. Hence
\begin{align}
 H(F\mid G)&\leq q^{d+1}\log_2q,\\
 H(X^{0:T-1}\mid G)&\leq
     H(X^0,F\mid G)
     \leq(n+q^{d+1})\log_2q.
 \label{eq:unknownruleentropy}
\end{align}
This avoids the invalid use of differential entropy as a bound on
discrete trajectory entropy.

For the binary pilot, $n=64$, $d+1=9$, $q=2$, $T=16$:
\begin{equation}
 H(X^{0:15}\mid G)\leq64+512=576\ {\rm bits}.
\end{equation}
The uncompressed 16-frame state array has $16\cdot64=1024$ binary entries.
With a known rule its entropy bound is only 64 bits. With an unknown
sampled rule the bound is at most 576 bits; the neural rule family and its
nonuniform prior can lower it further. These are distributional entropy
bounds, not a prediction of the finite gzip ratio.

For stochastic transitions, additional fresh randomness remains even
after conditioning on the parameters. The relevant decomposition becomes
\begin{equation}
 H(X^{0:T-1}\mid\omega,G)
 =H(X^0\mid\omega,G)+
   \sum_{t=0}^{T-2}H(X^{t+1}\mid X^t,\omega,G).
 \label{eq:stochasticentropy}
\end{equation}
The second term is generally nonzero. Thus the 576-bit argument must not
be applied to a categorical softmax sampler with nonzero temperature.

\subsection{Rule inference as an idealized prediction problem}

For context $C$ and next token $Y$, an ideal rule-aware predictor has
\begin{equation}
 p(Y\mid C,G)=
     \sum_f p(Y\mid C,f,G)\,p(f\mid C,G).
 \label{eq:posteriorpredictive}
\end{equation}
If there are unresolved initial states or stochastic transitions, the
conditional factor also integrates those uncertainties.
The transformer can approximate this computation without explicitly
storing a symbolic posterior; the equation alone does not establish that
it actually does so.

An exact toy calculation makes the distinction concrete. Suppose the
local binary rule is drawn uniformly from all $2^{512}$ truth tables,
independently of the initial state,
not from the pilot's Gaussian-weight MLP prior. An observed complete
trajectory queries local 9-bit patterns and records their next bits.
If $D(C)$ distinct patterns have been queried, consistency fixes precisely
those $D(C)$ entries. The remaining entries are still independent fair
bits, even though successive queries can depend on earlier answers:
\begin{equation}
 H(F\mid C=c,X^0=x^0,G)=512-D(c).
 \label{eq:truthposterior}
\end{equation}
For the random context, taking expectations gives
$I(F;C\mid X^0,G)=\E[D(C)]$ when $C$ contains only the resulting transition
observations with their queried patterns. Each new truth-table entry
provides one bit; repeated queries provide none.

For example, 100 distinct queried patterns leave $412$ unknown bits.
Predicting an already queried pattern has Bayes error zero. Predicting
a previously unqueried pattern still has Bayes error $1/2$ and loss
$\log2$ nats under this prior. Thus memorizing visited transitions does
not establish extrapolation to unvisited neighborhoods. New rules and
new initial conditions probe different forms of generalization.
The number $D(C)$ is not measured in the historical pilot.

\subsection{Contextual structure versus unigram matching}

For any predictor $q_\theta(Y\mid C)$ with adequate support,
\begin{equation}
 \E[-\log q_\theta(Y\mid C)]
 =H(Y\mid C)+
   \E_C\KL\!\left(p(\cdot\mid C)\,\|\,q_\theta(\cdot\mid C)\right),
 \label{eq:kl}
\end{equation}
using nats throughout this equation. An oracle unigram predictor attains
$H(Y)$. Its improvement to the oracle contextual predictor is exactly
\begin{equation}
 H(Y)-H(Y\mid C)=I(Y;C).
 \label{eq:mi}
\end{equation}
Matching token frequencies therefore matches $H(Y)$, not
$H(Y\mid C)$ or the difficulty of approximating the conditional model.
High marginal entropy alone is not evidence of useful temporal
computation.

The pilot shuffle is a uniform permutation of the patch-token multiset
within each trajectory, retaining frame markers. It does not produce
independent draws. For $M$ tokens with fraction $p$ of one chosen token,
indicators $A_j$ and $A_k$ at two different shuffled positions satisfy
\begin{align}
 \E[A_jA_k]&=\frac{(Mp)(Mp-1)}{M(M-1)},\\
 \Cov(A_j,A_k)&=-\frac{p(1-p)}{M-1}.
 \label{eq:shufflecov}
\end{align}
For $M=256$ and a hypothetical $p=1/2$ this is
$-1/1020=-0.0009803922$.
Markers also retain their deterministic positions. It is therefore
incorrect to set all contextual mutual information of the shuffled
control to zero.

\subsection{A task-focused loss}

For retained frames and patch tokens, define a mask $a_j$ that equals one
only for patch targets after the initial context frame, and zero for
markers or initial-frame targets. A useful synthetic objective is
\begin{equation}
 L_{\rm patch}(\theta)=
 -\E\!\left[\frac{1}{\sum_j a_j}
       \sum_j a_j\log p_\theta(Z_j\mid Z_{<j})\right].
 \label{eq:masked}
\end{equation}
For 16 retained frames and 16 patches/frame, this gives
$(16-1)16=240$ patch targets per complete trajectory.
Unmasked sequences have 32 marker tokens among 288 tokens, a fraction
$1/9$. Predictable marker positions can reduce total loss without proving
rule learning. The numerator, denominator and units must be specified
when comparing masked patch loss with unmasked next-token loss.

Nothing in \eqref{eq:mi} or \eqref{eq:masked} supplies factual information
about independently sampled language data. Any finite-budget advantage
must arise through the learned predictor/optimization bias. Whether this
bias is useful is the alignment question in Section~\ref{sec:transfer}.

